\section{Finite-Dimensional Vector Spaces}

\subsection{Span and Linear Independence}
\subsubsection{Linear Combinations and Span}
\begin{definition}[Linear Combination]
	A \textit{linear combination} of a list $v_1, \dots, v_m$ of vectors in $V$ is a vector of the form
	\[
	a_1v_1 + \cdots + a_mv_m,
	\]
	where $a_1, \dots, a_m \in \F$.
\end{definition}
\begin{definition}[Span]
	The set of all linear combinations of a list of vectors $v_1, \dots, v_m$ in $V$ is called the \textit{span} of $v_1, \dots, v_m$, denoted by $\text{span}(v_1, \dots, v_m)$. In other words,
	\[
	\text{span}(v_1, \dots, v_m) = \{a_1v_1 + \cdots + a_mv_m : a_1, \dots, a_m \in \F\}.
	\]
	The span of the empty list $( \, )$ is defined to be $\{0\}$.
\end{definition}
\begin{theorem}
	The span of a list of vectors in $\V$ is the smallest subspace of $\V$ containing all vectors in the list.
\end{theorem}
\begin{proof}
	Suppose $v_1, \dots, v_m$ is a list of vectors in $\V$.
	
	First we show that $\text{span}(v_1, \dots, v_m)$ is a subspace of $V$. The additive identity is in $\text{span}(v_1, \dots, v_m)$ because
	\[
	0 = 0v_1 + \cdots + 0v_m.
	\]
	Also, $\text{span}(v_1, \dots, v_m)$ is closed under addition because
	\[
	(a_1v_1 + \cdots + a_mv_m) + (c_1v_1 + \cdots + c_mv_m) = (a_1 + c_1)v_1 + \cdots + (a_m + c_m)v_m.
	\]
	Furthermore, $\text{span}(v_1, \dots, v_m)$ is closed under scalar multiplication because
	\[
	\lambda(a_1v_1 + \cdots + a_mv_m) = \lambda a_1v_1 + \cdots + \lambda a_mv_m.
	\]
	Thus $\text{span}(v_1, \dots, v_m)$ is a subspace of $\V$.
	
	Each $v_k$ is a linear combination of $v_1, \dots, v_m$ (to show this, set $a_k = 1$ and let the other $a$'s in equal 0). Thus $\text{span}(v_1, \dots, v_m)$ contains each $v_k$. Conversely, because subspaces are closed under scalar multiplication and addition, every subspace of $\V$ that contains each $v_k$ contains $\text{span}(v_1, \dots, v_m)$. Thus $\text{span}(v_1, \dots, v_m)$ is the smallest subspace of $\V$ containing all the vectors $v_1, \dots, v_m$.
\end{proof}
\begin{definition}
	A vector space is called \textit{finite-dimensional} if some list of vectors in it spans the space.
\end{definition}
\begin{definition}[Polynomial] \mbox{}
	\begin{itemize}
		\item A function $p: \F \to \F$ is called a \textit{polynomial} with coefficients in $\F$ if there exist $a_0, \dots, a_m \in \F$ such that
		\[
		p(z) = a_0 + a_1z + a_2z^2 + \cdots + a_mz^m
		\]
		for all $z \in \F$.
		\item $\mathcal{P}(\F)$ is the set of all polynomials with coefficients in $\F$.
	\end{itemize}
\end{definition}
\begin{definition}[Degree of a Polynomial] \mbox{}
	\begin{itemize}
		\item A polynomial $p \in \mathcal{P}(\F)$ is said to have \textit{degree} $m$ if there exist scalars $a_0, a_1, \dots, a_m \in \F$ with $a_m \neq 0$ such that for every $z \in \F$, we have
		\[
		p(z) = a_0 + a_1z + \cdots + a_mz^m.
		\]
		\item The polynomial that is identically $0$ is said to have degree $-\infty$. 
		\item The degree of a polynomial $p$ is denoted by $\deg p$.
	\end{itemize}
\end{definition}
\subsubsection{Linear Independence}
\begin{definition}[Linearly Independent] \mbox{}
	\begin{itemize}
		\item A list $v_1, \dots, v_m$ of vectors in $\V$ is called \textit{linearly independent} if the only choice of $a_1, \dots, a_m \in \F$ that makes
		\[
		a_1 v_1 + \dots + a_m v_m = 0
		\]
		is $a_1 = \dots = a_m = 0$.
		
		\item The empty list $( \ )$ is also declared to be linearly independent.
	\end{itemize}
\end{definition}
\begin{definition}[Linearly Dependent] \mbox{}
	\begin{itemize}
		\item A list of vectors in $\V$ is called \textit{linearly dependent} if it is not linearly independent.
		
		\item In other words, a list $v_1, \dots, v_m$ of vectors in $\V$ is linearly dependent if there exist $a_1, \dots, a_m \in \F$, not all $0$, such that $a_1 v_1 + \dots + a_m v_m = 0$.
	\end{itemize}
\end{definition}
\begin{theorem}[Linear Dependence Lemma]
	Suppose $v_1, \dots, v_m$ is a linearly dependent list in $\V$. Then there exists $k \in \{1, 2, \dots, m\}$ such that
	\[
	v_k \in \text{span}(v_1, \dots, v_{k-1}).
	\]
	Furthermore, if $k$ satisfies the condition above and the $k^{\text{th}}$ term is removed from $v_1, \dots, v_m$, then the span of the remaining list equals $\text{span}(v_1, \dots, v_m)$.
\end{theorem}
\begin{proof}
	Because the list $v_1, \dots, v_m$ is linearly dependent, there exist numbers $a_1, \dots, a_m \in \F$, not all $0$, such that
	\[
	a_1 v_1 + \dots + a_m v_m = 0.
	\]
	Let $k$ be the largest element of $\{1, \dots, m\}$ such that $a_k \neq 0$. Then
	\[
	v_k = -\frac{a_1}{a_k} v_1 - \dots - \frac{a_{k-1}}{a_k} v_{k-1},
	\]
	which proves that $v_k \in \text{span}(v_1, \dots, v_{k-1})$, as desired.
	
	Now suppose $k$ is any element of $\{1, \dots, m\}$ such that $v_k \in \text{span}(v_1, \dots, v_{k-1})$. Let $b_1, \dots, b_{k-1} \in \F$ be such that
	\begin{equation} \label{eq:1}
		v_k = b_1 v_1 + \dots + b_{k-1} v_{k-1}.
	\end{equation}
	Suppose $u \in \text{span}(v_1, \dots, v_m)$. Then there exist $c_1, \dots, c_m \in \F$ such that
	\[
	u = c_1 v_1 + \dots + c_m v_m.
	\]
	In the equation above, we can replace $v_k$ with the right side of \eqref{eq:1}, which shows that $u$ is in the span of the list obtained by removing the $k^{\text{th}}$ term from $v_1, \dots, v_m$. Thus removing the $k^{\text{th}}$ term of the list $v_1, \dots, v_m$ does not change the span of the list.
\end{proof}
\begin{theorem}
	In a finite-dimensional vector space, the length of every linearly independent list of vectors is less than or equal to the length of every spanning list of vectors.
\end{theorem}
\begin{proof}
	Suppose that $u_1, \dots, u_m$ is linearly independent in $\V$. Suppose also that $w_1, \dots, w_n$ spans $\V$. We need to prove that $m \le n$. We do so through the process described below with $m$ steps; note that in each step we add one of the $u$'s and remove one of the $w$'s.
	
	\noindent \textbf{Step 1} \\
	Let $B$ be the list $w_1, \dots, w_n$, which spans $\V$. Adjoining $u_1$ at the beginning of this list produces a linearly dependent list (because $u_1$ can be written as a linear combination of $w_1, \dots, w_n$). In other words, the list
	\[
	u_1, w_1, \dots, w_n
	\]
	is linearly dependent.
	
	Thus by the linear dependence lemma, one of the vectors in the list above is a linear combination of the previous vectors in the list. We know that $u_1 \neq 0$ because the list $u_1, \dots, u_m$ is linearly independent. Thus $u_1$ is not in the span of the previous vectors in the list above (because $u_1$ is not in $\{0\}$, which is the span of the empty list). Hence the linear dependence lemma implies that we can remove one of the $w$'s so that the new list $B$ (of length $n$) consisting of $u_1$ and the remaining $w$'s spans $\V$.
	
	\noindent \textbf{Step $k$, for $k = 2, \dots, m$} \\
	The list $B$ (of length $n$) from step $k - 1$ spans $\V$. In particular, $u_k$ is in the span of the list $B$. Thus the list of length $(n + 1)$ obtained by adjoining $u_k$ to $B$, placing it just after $u_1, \dots, u_{k-1}$, is linearly dependent. By the linear dependence lemma (2.9), one of the vectors in this list is in the span of the previous ones, and because $u_1, \dots, u_k$ is linearly independent, this vector cannot be one of the $u$'s.
	
	Hence there still must be at least one remaining $w$ at this step. We can remove from our new list (after adjoining $u_k$ in the proper place) a $w$ that is a linear combination of the previous vectors in the list, so that the new list $B$ (of length $n$) consisting of $u_1, \dots, u_k$ and the remaining $w$'s spans $\V$.
	
	After step $m$, we have added all the $u$'s and the process stops. At each step as we add a $u$ to $B$, the linear dependence lemma implies that there is some $w$ to remove. Thus there are at least as many $w$'s as $u$'s.
\end{proof}
\begin{theorem}
	Every subspace of a finite-dimensional vector space is finite-dimensional.
\end{theorem}
\begin{proof}
	Suppose $\V$ is finite-dimensional and $U$ is a subspace of $\V$. We need to prove that $U$ is finite-dimensional. We do this through the following multistep construction.
	
	\noindent \textbf{Step 1} \\
	If $U = \{0\}$, then $U$ is finite-dimensional and we are done. If $U \neq \{0\}$, then choose a nonzero vector $u_1 \in U$.
	
	\noindent \textbf{Step $k$} \\
	If $U = \text{span}(u_1, \dots, u_{k-1})$, then $U$ is finite-dimensional and we are done. If $U \neq \text{span}(u_1, \dots, u_{k-1})$, then choose a vector $u_k \in U$ such that
	\[
	u_k \notin \text{span}(u_1, \dots, u_{k-1}).
	\]
	
	After each step, as long as the process continues, we have constructed a list of vectors such that no vector in this list is in the span of the previous vectors. Thus after each step we have constructed a linearly independent list, by the linear dependence lemma. This linearly independent list cannot be longer than any spanning list of $\V$. Thus the process eventually terminates, which means that $U$ is finite-dimensional.
\end{proof}
\subsubsection{Exercise 2A}
\setcounterpageref{question}{0}
\begin{question}
	Find a list of four distinct vectors in $\F^3$ whose span equals
	\[
	\{(x, y, z) \in \F^3 : x + y + z = 0\}.
	\]
\end{question}
\begin{proof}
	Let $U = \{(x, y, z) \in \F^3 : x + y + z = 0\}$. Notice that any vector $(x, y, z)$ in $U$ must satisfy $z = -x - y$. Therefore, we can express any vector in $U$ as:
	\[
	(x, y, -x-y) = x(1, 0, -1) + y(0, 1, -1)
	\]
	This implies that the list of two vectors $(1, 0, -1)$ and $(0, 1, -1)$ spans $U$. 
	
	To obtain a list of four distinct vectors that span $U$, we can simply include these two vectors and add any two other distinct vectors that also lie in $U$. We can generate other vectors in $U$ by choosing arbitrary values for $x$ and $y$ and setting $z = -x - y$:
	\begin{itemize}
		\item Let $x=1, y=0 \implies v_1 = (1, 0, -1)$
		\item Let $x=0, y=1 \implies v_2 = (0, 1, -1)$
		\item Let $x=1, y=1 \implies v_3 = (1, 1, -2)$
		\item Let $x=-1, y=0 \implies v_4 = (-1, 0, 1)$
	\end{itemize}
	
	Since $v_1$ and $v_2$ are in our list, the span of $(v_1, v_2, v_3, v_4)$ contains all linear combinations of $v_1$ and $v_2$, which means it contains all of $U$. Furthermore, since $v_3, v_4 \in U$, their inclusion does not expand the span beyond $U$. 
	
	Thus, the span of the list is exactly $U$, and a valid list of four distinct vectors is:
	\[
	(1, 0, -1), \ (0, 1, -1), \ (1, 1, -2), \ (-1, 0, 1)
	\]
\end{proof}
\begin{question}
	Prove or give a counterexample: If $v_1, v_2, v_3, v_4$ spans $\V$, then the list
	\[
	v_1 - v_2, v_2 - v_3, v_3 - v_4, v_4
	\]
	also spans $\V$.
\end{question}
\begin{proof}
	The statement is true. 
	
	\textbf{Proof:}
	Let $u \in \V$. Since the list $v_1, v_2, v_3, v_4$ spans $\V$, there exist scalars $a_1, a_2, a_3, a_4 \in \F$ such that:
	\[
	u = a_1 v_1 + a_2 v_2 + a_3 v_3 + a_4 v_4
	\]
	We must show that $u$ can be expressed as a linear combination of the vectors in the new list: $v_1 - v_2, v_2 - v_3, v_3 - v_4, v_4$. 
	
	Notice that we can rewrite each of the original vectors in terms of the new vectors by using telescoping sums:
	\begin{align*}
		v_4 &= v_4 \\
		v_3 &= (v_3 - v_4) + v_4 \\
		v_2 &= (v_2 - v_3) + (v_3 - v_4) + v_4 \\
		v_1 &= (v_1 - v_2) + (v_2 - v_3) + (v_3 - v_4) + v_4
	\end{align*}
	
	Substitute these expressions back into the original equation for $u$:
	\begin{align*}
		u &= a_1\big((v_1 - v_2) + (v_2 - v_3) + (v_3 - v_4) + v_4\big) \\
		&\quad + a_2\big((v_2 - v_3) + (v_3 - v_4) + v_4\big) \\
		&\quad + a_3\big((v_3 - v_4) + v_4\big) \\
		&\quad + a_4v_4
	\end{align*}
	
	Distributing the scalars and factoring out the new vectors yields:
	\[
	u = a_1(v_1 - v_2) + (a_1 + a_2)(v_2 - v_3) + (a_1 + a_2 + a_3)(v_3 - v_4) + (a_1 + a_2 + a_3 + a_4)v_4
	\]
	
	Because the field of scalars $\F$ is closed under addition, the coefficients $a_1$, $(a_1 + a_2)$, $(a_1 + a_2 + a_3)$, and $(a_1 + a_2 + a_3 + a_4)$ are all valid scalars in $\F$. 
	
	Thus, any arbitrary vector $u \in \V$ can be written as a linear combination of $v_1 - v_2, v_2 - v_3, v_3 - v_4, v_4$, meaning this list also spans $\V$.
\end{proof}
\begin{question}
	Suppose $v_1, \dots, v_m$ is a list of vectors in $\V$. For $k \in \{1, \dots, m\}$, let
	\[
	w_k = v_1 + \dots + v_k.
	\]
	Show that $\operatorname{span}(v_1, \dots, v_m) = \operatorname{span}(w_1, \dots, w_m)$.
\end{question}
\begin{proof}
	To show that $\operatorname{span}(v_1, \dots, v_m) = \operatorname{span}(w_1, \dots, w_m)$, we must prove that each set is a subset of the other.
	
	\textbf{Part 1: $\operatorname{span}(w_1, \dots, w_m) \subseteq \operatorname{span}(v_1, \dots, v_m)$}
	
	By definition, for each $k \in \{1, \dots, m\}$, the vector $w_k$ is given by:
	\[
	w_k = v_1 + \dots + v_k
	\]
	This shows that every $w_k$ is a linear combination of the vectors in the list $v_1, \dots, v_m$. Since every $w_k \in \operatorname{span}(v_1, \dots, v_m)$, and the span is a vector space closed under addition and scalar multiplication, any linear combination of $w_1, \dots, w_m$ must also belong to $\operatorname{span}(v_1, \dots, v_m)$. Therefore, 
	\[
	\operatorname{span}(w_1, \dots, w_m) \subseteq \operatorname{span}(v_1, \dots, v_m).
	\]
	
	\textbf{Part 2: $\operatorname{span}(v_1, \dots, v_m) \subseteq \operatorname{span}(w_1, \dots, w_m)$}
	
	We can express each vector $v_k$ as a linear combination of the $w$ vectors. For $k = 1$, we clearly have:
	\[
	v_1 = w_1
	\]
	For any $k \in \{2, \dots, m\}$, we can isolate $v_k$ by subtracting consecutive $w$ vectors:
	\begin{align*}
		w_k - w_{k-1} &= (v_1 + \dots + v_{k-1} + v_k) - (v_1 + \dots + v_{k-1}) \\
		&= v_k
	\end{align*}
	This demonstrates that every $v_k$ can be written as a linear combination of $w_1, \dots, w_m$ (specifically, $w_k - w_{k-1}$). Since every $v_k \in \operatorname{span}(w_1, \dots, w_m)$, it follows that any linear combination of $v_1, \dots, v_m$ is also in $\operatorname{span}(w_1, \dots, w_m)$. Therefore,
	\[
	\operatorname{span}(v_1, \dots, v_m) \subseteq \operatorname{span}(w_1, \dots, w_m).
	\]
	
	 $\operatorname{span}(w_1, \dots, w_m) \subseteq \operatorname{span}(v_1, \dots, v_m)$ and $\operatorname{span}(v_1, \dots, v_m) \subseteq \operatorname{span}(w_1, \dots, w_m)$, the two sets must be equal.
\end{proof}
\begin{question} \mbox{}
	\begin{enumerate}
	\item[(a)] Show that a list of length one in a vector space is linearly independent if and only if the vector in the list is not $0$.
	\item[(b)] Show that a list of length two in a vector space is linearly independent if and only if neither of the two vectors in the list is a scalar multiple of the other.
	\end{enumerate}
\end{question}
\begin{proof} \mbox{}\\
	\textbf{Part (a)} \mbox{} \\
	($\implies$) Suppose the list containing only $v$ is linearly independent. By definition of linear independence, the equation $av = 0$ implies $a = 0$ for $a \in \F$. If $v$ were equal to the zero vector ($v = 0$), then any non-zero scalar $a \in \F$ (such as $a = 1$) would satisfy $1 \cdot 0 = 0$. This would mean a non-trivial linear combination equals zero, making the list linearly dependent. This is a contradiction. Therefore, $v \neq 0$.
	
	($\impliedby$) Suppose $v \neq 0$. Let $a \in \F$ be a scalar such that $av = 0$. Since $v$ is not the zero vector, it must be that $a = 0$ (otherwise, we could multiply both sides by the multiplicative inverse $1/a$ to get $v = (1/a)0 = 0$, contradicting $v \neq 0$). Since $a = 0$ is the only solution to $av = 0$, the list of length one is linearly independent.
	
	\vspace{1em}
	
	\textbf{Part (b)}
	
	Let $v_1, v_2$ be a list of length two in a vector space $V$. 
	
	($\implies$) Suppose the list $v_1, v_2$ is linearly independent. We must show that neither vector is a scalar multiple of the other. Assume, for the sake of contradiction, that one is a scalar multiple of the other. 
	If $v_1 = c v_2$ for some scalar $c \in \F$, then we can write:
	\[
	1 v_1 - c v_2 = 0
	\]
	Since the coefficient of $v_1$ is $1 \neq 0$, this is a non-trivial linear combination that equals $0$, which contradicts the assumption of linear independence. 
	Similarly, if $v_2 = c v_1$, then $-c v_1 + 1 v_2 = 0$, which again contradicts linear independence since the coefficient of $v_2$ is $1 \neq 0$. Thus, neither vector can be a scalar multiple of the other.
	
	($\impliedby$) Suppose neither vector is a scalar multiple of the other. We must show that the list $v_1, v_2$ is linearly independent. Suppose there exist scalars $a_1, a_2 \in \F$ such that:
	\[
	a_1 v_1 + a_2 v_2 = 0
	\]
	If $a_1 \neq 0$, we can rearrange the equation and divide by $a_1$ to isolate $v_1$:
	\[
	v_1 = -\frac{a_2}{a_1} v_2
	\]
	This implies that $v_1$ is a scalar multiple of $v_2$, which contradicts our initial assumption. Therefore, we must have $a_1 = 0$. 
	
	Substituting $a_1 = 0$ into the original equation gives:
	\[
	a_2 v_2 = 0
	\]
	If $a_2 \neq 0$, then $v_2 = 0$. However, if $v_2 = 0$, then $v_2 = 0 v_1$, which means $v_2$ is a scalar multiple of $v_1$, again contradicting our assumption. Therefore, $a_2$ must also be $0$. 
	
	Since the only solution is $a_1 = 0$ and $a_2 = 0$, the list $v_1, v_2$ is linearly independent.
\end{proof}
\begin{question}
	Find a number $t$ such that
	\[
	(3, 1, 4), (2, -3, 5), (5, 9, t)
	\]
	is not linearly independent in $\R^3$.
\end{question}
\begin{proof}
	We can set up the equation:
	\[
	(5, 9, t) = a(3, 1, 4) + b(2, -3, 5)
	\]
	This yields a system of three linear equations:
	\begin{align}
		3a + 2b &= 5 \\
		a - 3b &= 9 \\
		4a + 5b &= t
	\end{align}
	Thus, the number is $t = 2$.
\end{proof}
\begin{question}
	Show that the list $(2, 3, 1), (1, -1, 2), (7, 3, c)$ is linearly dependent in $\F^3$ if and only if $c = 8$.
\end{question}
\begin{proof}
	For the list of three vectors to be linearly dependent, the third vector must be a linear combination of the first two. 
	
	This means we can set up the equation:
	\[
	a(2, 3, 1) + b(1, -1, 2) = (7, 3, c)
	\]
	This produces a system of three linear equations:
	\begin{align}
		2a + b &= 7 \\
		3a - b &= 3 \\
		a + 2b &= c
	\end{align}
	
	We can solve for $a$ and $b$ using the first two equations. Adding equation (1) and equation (2) together yields:
	\begin{align*}
		a &= 2
		b &= 3
		c &= 8
	\end{align*}
	
	Because the system is only consistent when $c = 8$, the third vector is a linear combination of the first two if and only if $c = 8$. Thus, the list is linearly dependent if and only if $c = 8$.
\end{proof}
\begin{question} \mbox{}\\
	\begin{enumerate} 
		\item[(a)] Show that if we think of $\C$ as a vector space over $\R$, then the list $1 + i, 1 - i$ is linearly independent.
		\item[(b)] Show that if we think of $\C$ as a vector space over $\C$, then the list $1 + i, 1 - i$ is linearly dependent.
	\end{enumerate}
\end{question}
\begin{proof}
	\textbf{Part (a)}
	We must show that the only real scalars $a, b \in \R$ that satisfy the equation
	\[
	a(1 + i) + b(1 - i) = 0
	\]
	are $a = 0$ and $b = 0$. 
	
	Expanding the equation, we get:
	\[
	a + ai + b - bi = 0
	\]
	Grouping the real and imaginary parts gives:
	\[
	(a + b) + (a - b)i = 0
	\]
	For a complex number to equal zero, both its real and imaginary parts must be zero. This yields the system of linear equations:
	\begin{align*}
		a + b &= 0 \\
		a - b &= 0
	\end{align*}
	Which implies $a = 0$ and $b = 0$. 
	
	Since the only solution is $a = 0$ and $b = 0$, the list $1 + i, 1 - i$ is linearly independent over $\R$.
	
	\vspace{1em}
	
	\textbf{Part (b)}
	
	Now suppose $\C$ is a vector space over $\C$.
	
	This allows us to write:
	\[
	1(1 + i) - i(1 - i) = 0
	\]
	Here, our chosen scalars are $c_1 = 1$ and $c_2 = -i$. Because $c_1 \neq 0$ (and $c_2 \neq 0$), we have a non-trivial linear combination that equals zero. 
	
	Therefore, the list $1 + i, 1 - i$ is linearly dependent over $\C$. 
\end{proof}
\begin{question}
	Suppose $v_1, v_2, v_3, v_4$ is linearly independent in $\V$. Prove that the list
	\[
	v_1 - v_2, v_2 - v_3, v_3 - v_4, v_4
	\]
	is also linearly independent.
\end{question}
\begin{proof}
	To prove that the list $v_1 - v_2, v_2 - v_3, v_3 - v_4, v_4$ is linearly independent, suppose there exist scalars $a_1, a_2, a_3, a_4 \in \F$ such that:
	\[
	a_1(v_1 - v_2) + a_2(v_2 - v_3) + a_3(v_3 - v_4) + a_4v_4 = 0
	\]
	
	Expanding the terms and regrouping by the vectors $v_1, v_2, v_3, v_4$ gives:
	\begin{align*}
		a_1v_1 - a_1v_2 + a_2v_2 - a_2v_3 + a_3v_3 - a_3v_4 + a_4v_4 &= 0 \\
		a_1v_1 + (a_2 - a_1)v_2 + (a_3 - a_2)v_3 + (a_4 - a_3)v_4 &= 0
	\end{align*}
	
	We are given that the list $v_1, v_2, v_3, v_4$ is linearly independent. By the definition of linear independence, the only solution to this equation is when all coefficients are equal to zero. This gives us the following system of equations:
	\begin{align*}
		a_1 &= 0 \\
		a_2 - a_1 &= 0 \\
		a_3 - a_2 &= 0 \\
		a_4 - a_3 &= 0
	\end{align*}
	Because the only solution to the linear combination is the trivial solution $a_1 = a_2 = a_3 = a_4 = 0$, the list $v_1 - v_2, v_2 - v_3, v_3 - v_4, v_4$ is also linearly independent.
\end{proof}
\begin{question}
	Prove or give a counterexample: If $v_1, v_2, \dots, v_m$ is a linearly independent list of vectors in $\V$, then
	\[
	5v_1 - 4v_2, v_2, v_3, \dots, v_m
	\]
	is linearly independent.
\end{question}
\begin{proof}
	The statement is true.
	
	Suppose there exist scalars $a_1, a_2, \dots, a_m \in \F$ such that:
	\[
	a_1(5v_1 - 4v_2) + a_2v_2 + a_3v_3 + \dots + a_mv_m = 0
	\]
	
	Expanding the terms and regrouping by the vectors $v_1, v_2, \dots, v_m$ gives:
	\begin{align*}
		5a_1v_1 - 4a_1v_2 + a_2v_2 + a_3v_3 + \dots + a_mv_m &= 0 \\
		5a_1v_1 + (a_2 - 4a_1)v_2 + a_3v_3 + \dots + a_mv_m &= 0
	\end{align*}
	
	Because we are given that the list $v_1, v_2, \dots, v_m$ is linearly independent, the coefficients of each vector in this linear combination must be exactly zero. This produces the following system of equations:
	\begin{align*}
		5a_1 &= 0 \\
		a_2 - 4a_1 &= 0 \\
		a_3 &= 0 \\
		&\vdots \\
		a_m &= 0
	\end{align*}
	
	From the first equation, dividing by $5$ yields $a_1 = 0$. 
	Substituting $a_1 = 0$ into the second equation yields $a_2 - 4(0) = 0$, which gives $a_2 = 0$. 
	The remaining equations directly state that $a_3 = 0, \dots, a_m = 0$.
	
	Since the only solution is the trivial solution $a_1 = a_2 = \dots = a_m = 0$, the list $5v_1 - 4v_2, v_2, v_3, \dots, v_m$ is linearly independent.
\end{proof}
\begin{question}
	Prove or give a counterexample: If $v_1, v_2, \dots, v_m$ is a linearly independent list of vectors in $\V$ and $\lambda \in \F$ with $\lambda \neq 0$, then $\lambda v_1, \lambda v_2, \dots, \lambda v_m$ is linearly independent.
\end{question}
\begin{proof}
	The statement is true.
	
	Suppose there exist scalars $a_1, a_2, \dots, a_m \in \F$ such that:
	\[
	a_1(\lambda v_1) + a_2(\lambda v_2) + \dots + a_m(\lambda v_m) = 0
	\]
	
	Using the associativity of scalar multiplication in a vector space, we can rewrite this equation as:
	\[
	(\lambda a_1)v_1 + (\lambda a_2)v_2 + \dots + (\lambda a_m)v_m = 0
	\]
	
	We are given that the list $v_1, v_2, \dots, v_m$ is linearly independent. By the definition of linear independence, the scalar coefficients of the vectors in this linear combination must all equal zero. This gives:
	\[
	\lambda a_1 = 0, \quad \lambda a_2 = 0, \quad \dots, \quad \lambda a_m = 0
	\]
	
	We are also given that $\lambda \neq 0$. Because $\lambda$ is a non-zero element of the field $\mathbf{F}$, it has a multiplicative inverse, $\lambda^{-1}$. Multiplying both sides of each equation by $\lambda^{-1}$ yields:
	\[
	a_1 = 0, \quad a_2 = 0, \quad \dots, \quad a_m = 0
	\]
	
	Because the only scalars that satisfy the initial equation are $a_1 = a_2 = \dots = a_m = 0$, the list $\lambda v_1, \lambda v_2, \dots, \lambda v_m$ is linearly independent.
\end{proof}
\begin{question}
	Prove or give a counterexample: If $v_1, \dots, v_m$ and $w_1, \dots, w_m$ are linearly independent lists of vectors in $\V$, then the list $v_1 + w_1, \dots, v_m + w_m$ is linearly independent.
\end{question}
\begin{proof}
		The statement is false. 
	
	Define $w_i = -v_i$ for each $i = 1, \dots, m$.
	
	First, we verify that $w_1, \dots, w_m$ is linearly independent. Suppose $a_1, \dots, a_m \in \mathbb{F}$ are scalars such that:
	\[
	a_1 w_1 + \dots + a_m w_m = 0
	\]
	Substituting $w_i = -v_i$, we obtain:
	\[
	-a_1 v_1 - \dots - a_m v_m = 0
	\]
	Since $v_1, \dots, v_m$ is linearly independent, all coefficients must be zero: $-a_1 = 0, \dots, -a_m = 0$, which means $a_1 = 0, \dots, a_m = 0$. Thus, $w_1, \dots, w_m$ is a linearly independent list.
	
	Now, consider the list $v_1 + w_1, \dots, v_m + w_m$. For each $i$, we have:
	\[
	v_i + w_i = v_i - v_i = 0
	\]
	The new list is exactly $0, \dots, 0$. Because this list contains the zero vector, it is linearly dependent. This disproves the statement.
\end{proof}
\begin{question}
	Suppose $v_1, \dots, v_m$ is linearly independent in $\V$ and $w \in \V$. Prove that if $v_1 + w, \dots, v_m + w$ is linearly dependent, then $w \in \text{span}(v_1, \dots, v_m)$.
\end{question}
\begin{proof}
	Because $v_1 + w, \dots, v_m + w$ is linearly dependent, there exist scalars $a_1, \dots, a_m \in \mathbb{F}$, not all zero, such that
	\[
	a_1(v_1 + w) + \dots + a_m(v_m + w) = 0.
	\]
	Expanding and rearranging the terms, we obtain
	\[
	a_1 v_1 + \dots + a_m v_m + (a_1 + \dots + a_m)w = 0.
	\]
	Let $c = a_1 + \dots + a_m$. If $c = 0$, then the equation simplifies to
	\[
	a_1 v_1 + \dots + a_m v_m = 0.
	\]
	Since $v_1, \dots, v_m$ is linearly independent, this implies that $a_1 = 0, \dots, a_m = 0$. However, this contradicts the fact that the scalars $a_1, \dots, a_m$ are not all zero. Therefore, it must be true that $c \neq 0$. 
	
	Since $c \neq 0$, we can divide by $c$ and solve for $w$:
	\[
	w = -\frac{a_1}{c}v_1 - \dots - \frac{a_m}{c}v_m.
	\]
	This demonstrates that $w$ is a linear combination of $v_1, \dots, v_m$. Thus, $w \in \text{span}(v_1, \dots, v_m)$.
\end{proof}
\begin{question}
	Suppose $v_1, \dots, v_m$ is linearly independent in $\V$ and $w \in \V$. Show that
	\[
	v_1, \dots, v_m, w \text{ is linearly independent} \iff w \notin \text{span}(v_1, \dots, v_m).
	\]
\end{question}
\begin{proof}
	First, we prove the forward direction ($\implies$). Suppose $v_1, \dots, v_m, w$ is linearly independent. Assume, for the sake of contradiction, that $w \in \text{span}(v_1, \dots, v_m)$. Then there exist scalars $a_1, \dots, a_m \in \mathbb{F}$ such that
	\[
	w = a_1 v_1 + \dots + a_m v_m.
	\]
	Rearranging this equation gives
	\[
	a_1 v_1 + \dots + a_m v_m - w = 0.
	\]
	This is a linear combination of $v_1, \dots, v_m, w$ that equals $0$, where the coefficients are not all zero (specifically, the coefficient of $w$ is $-1$). This contradicts the assumption that $v_1, \dots, v_m, w$ is linearly independent. Therefore, it must be true that $w \notin \text{span}(v_1, \dots, v_m)$.
	
	Next, we prove the backward direction ($\impliedby$). Suppose $w \notin \text{span}(v_1, \dots, v_m)$. To show that $v_1, \dots, v_m, w$ is linearly independent, suppose there exist scalars $a_1, \dots, a_m, b \in \mathbb{F}$ such that
	\[
	a_1 v_1 + \dots + a_m v_m + b w = 0.
	\]
	If $b \neq 0$, we can solve for $w$:
	\[
	w = -\frac{a_1}{b} v_1 - \dots - \frac{a_m}{b} v_m.
	\]
	This would imply $w \in \text{span}(v_1, \dots, v_m)$, which contradicts our assumption. Therefore, we must have $b = 0$. Substituting $b = 0$ into our equation yields
	\[
	a_1 v_1 + \dots + a_m v_m = 0.
	\]
	Because $v_1, \dots, v_m$ is linearly independent, we conclude that $a_1 = 0, \dots, a_m = 0$. Since all scalars $a_1, \dots, a_m, b$ are zero, the list $v_1, \dots, v_m, w$ is linearly independent.
\end{proof}
\begin{question}
	Suppose $v_1, \dots, v_m$ is a list of vectors in $V$. For $k \in \{1, \dots, m\}$, let
	\[
	w_k = v_1 + \dots + v_k.
	\]
	Show that the list $v_1, \dots, v_m$ is linearly independent if and only if the list $w_1, \dots, w_m$ is linearly independent.
\end{question}
\begin{proof}
	First, we prove the forward direction ($\implies$). Suppose $v_1, \dots, v_m$ is linearly independent. To show that $w_1, \dots, w_m$ is linearly independent, suppose there exist scalars $a_1, \dots, a_m \in \mathbb{F}$ such that
	\[
	a_1 w_1 + a_2 w_2 + \dots + a_m w_m = 0.
	\]
	Substituting $w_k = v_1 + \dots + v_k$ for each $k$, we obtain
	\[
	a_1(v_1) + a_2(v_1 + v_2) + \dots + a_m(v_1 + \dots + v_m) = 0.
	\]
	Rearranging the terms by factoring out each $v_k$, we get
	\[
	(a_1 + a_2 + \dots + a_m)v_1 + (a_2 + \dots + a_m)v_2 + \dots + a_m v_m = 0.
	\]
	Because $v_1, \dots, v_m$ is linearly independent, all coefficients must be zero:
	\begin{align*}
		a_1 + a_2 + \dots + a_m &= 0 \\
		a_2 + \dots + a_m &= 0 \\
		&\vdots \\
		a_m &= 0
	\end{align*}
	Working from the bottom up, $a_m = 0$, which implies $a_{m-1} = 0$, and so on, until $a_1 = 0$. Since all scalars $a_1, \dots, a_m$ are zero, $w_1, \dots, w_m$ is linearly independent.
	
	Next, we prove the backward direction ($\impliedby$). Suppose $w_1, \dots, w_m$ is linearly independent. To show that $v_1, \dots, v_m$ is linearly independent, suppose there exist scalars $b_1, \dots, b_m \in \mathbb{F}$ such that
	\[
	b_1 v_1 + b_2 v_2 + \dots + b_m v_m = 0.
	\]
	Notice that $v_1 = w_1$, and for $k \ge 2$, $v_k = w_k - w_{k-1}$. Substituting these expressions into our equation yields
	\[
	b_1(w_1) + b_2(w_2 - w_1) + b_3(w_3 - w_2) + \dots + b_m(w_m - w_{m-1}) = 0.
	\]
	Rearranging the terms to factor out each $w_k$, we obtain
	\[
	(b_1 - b_2)w_1 + (b_2 - b_3)w_2 + \dots + (b_{m-1} - b_m)w_{m-1} + b_m w_m = 0.
	\]
	Because $w_1, \dots, w_m$ is linearly independent, all coefficients must be zero:
	\begin{align*}
		b_1 - b_2 &= 0 \\
		b_2 - b_3 &= 0 \\
		&\vdots \\
		b_{m-1} - b_m &= 0 \\
		b_m &= 0
	\end{align*}
	From $b_m = 0$, it follows that $b_{m-1} = 0$, and continuing this chain gives $b_1 = b_2 = \dots = b_m = 0$. Therefore, $v_1, \dots, v_m$ is linearly independent.
\end{proof}
\begin{question}
	Explain why there does not exist a list of six polynomials that is linearly independent in $\mathcal{P}_4(\F)$.
\end{question}
\begin{proof}
	The vector space $\mathcal{P}_4(\F)$ consists of all polynomials with coefficients in $\F$ of degree at most 4. This space is spanned by the list of polynomials $1, z, z^2, z^3, z^4$.
	
	The length of this spanning list is 5. In any finite-dimensional vector space, the length of every linearly independent list of vectors is less than or equal to the length of every spanning list of vectors. 
	
	Since $\mathcal{P}_4(\F)$ has a spanning list of length 5, any linearly independent list in $\mathcal{P}_4(\F)$ must have a length of at most 5. Therefore, a list of six polynomials in $\mathcal{P}_4(\F)$ cannot be linearly independent.
\end{proof}
\begin{question}
	Explain why no list of four polynomials spans $\mathcal{P}_4(\F)$.
\end{question}
\begin{proof}
	The list of polynomials $1, z, z^2, z^3, z^4$ is linearly independent in $\mathcal{P}_4(\F)$. The length of this linearly independent list is 5. 
	
	In any finite-dimensional vector space, the length of every spanning list of vectors is greater than or equal to the length of every linearly independent list of vectors. 
	
	Because $\mathcal{P}_4(\F)$ contains a linearly independent list of length 5, any spanning list of $\mathcal{P}_4(\F)$ must have a length of at least 5. Therefore, no list of four polynomials can span the space.
\end{proof}
\begin{question}
	Prove that $\V$ is infinite-dimensional if and only if there is a sequence $v_1, v_2, \dots$ of vectors in $\V$ such that $v_1, \dots, v_m$ is linearly independent for every positive integer $m$.
\end{question}
\begin{proof}
	First, we prove the forward direction ($\implies$). Suppose $\V$ is infinite-dimensional. We will construct the sequence $v_1, v_2, \dots$ recursively. 
	Since $\V$ is infinite-dimensional, $\V \neq \{0\}$, so there exists a nonzero vector $v_1 \in \V$. The list $v_1$ is linearly independent. 
	Now, suppose for some positive integer $m$, we have chosen $v_1, \dots, v_m$ such that the list is linearly independent. Because $\V$ is infinite-dimensional, $v_1, \dots, v_m$ cannot span $\V$ (if it did, $\V$ would be finite-dimensional). Therefore, there exists a vector $v_{m+1} \in \V$ such that $v_{m+1} \notin \text{span}(v_1, \dots, v_m)$. Because $v_{m+1}$ is not in the span of the linearly independent list $v_1, \dots, v_m$, the extended list $v_1, \dots, v_m, v_{m+1}$ is also linearly independent. By mathematical induction, this process can be continued indefinitely to produce a sequence $v_1, v_2, \dots$ such that every finite prefix $v_1, \dots, v_m$ is linearly independent.
	
	Next, we prove the backward direction ($\impliedby$). Suppose there exists a sequence $v_1, v_2, \dots$ of vectors in $\V$ such that $v_1, \dots, v_m$ is linearly independent for every positive integer $m$. Assume, for the sake of contradiction, that $\V$ is finite-dimensional. Then $\V$ is spanned by some list of finite length, let's call this length $n$. 
	The length of any linearly independent list in a finite-dimensional vector space is less than or equal to the length of every spanning list. By our assumption, the list $v_1, \dots, v_{n+1}$ is linearly independent. The length of this linearly independent list is $n+1$, which is strictly greater than the length of the spanning list $n$. This is a contradiction. Therefore, our assumption that $\V$ is finite-dimensional must be false, meaning $\V$ is infinite-dimensional.
\end{proof}
\begin{question}
	Prove that $\F^\infty$ is infinite-dimensional.
\end{question}
\begin{proof}
	For each positive integer $k$, let $e_k \in \F^\infty$ be the sequence with $1$ in the $k$-th position and $0$ in all other positions. 
	
	Consider an arbitrary positive integer $m$. We will show that the list $e_1, \dots, e_m$ is linearly independent. Suppose there exist scalars $a_1, \dots, a_m \in \F$ such that
	\[
	a_1 e_1 + \dots + a_m e_m = 0.
	\]
	The left side of this equation evaluates to the sequence $(a_1, \dots, a_m, 0, 0, \dots)$, while the right side is the zero sequence $(0, 0, 0, \dots)$. Equating the components, we immediately obtain $a_1 = 0, \dots, a_m = 0$. Thus, the list $e_1, \dots, e_m$ is linearly independent.
	
	Because $\F^\infty$ contains a linearly independent list of length $m$ for every positive integer $m$, it cannot be spanned by any finite list of vectors (the length of a linearly independent list cannot exceed the length of any spanning list). Therefore, $\F^\infty$ is infinite-dimensional.
\end{proof}
\begin{question}
	Prove that the real vector space of all continuous real-valued functions on the interval $[0, 1]$ is infinite-dimensional.
\end{question}
\begin{proof}
	Let $\V$ be the real vector space of all continuous real-valued functions on the interval $[0, 1]$. Consider the functions $f_k \in \V$ defined by $f_k(x) = x^k$ for $k = 0, 1, 2, \dots$. 
	
	For any positive integer $m$, consider the list of $m$ functions $f_0, f_1, \dots, f_{m-1}$. Suppose there exist real scalars $a_0, a_1, \dots, a_{m-1}$ such that
	\[
	a_0 f_0 + a_1 f_1 + \dots + a_{m-1} f_{m-1} = 0.
	\]
	This implies that for all $x \in [0, 1]$, we have
	\[
	a_0 + a_1 x + \dots + a_{m-1} x^{m-1} = 0.
	\]
	The left side of this equation is a polynomial of degree at most $m-1$. A non-zero polynomial of degree at most $m-1$ can have at most $m-1$ roots. However, this polynomial equals zero for all infinitely many values of $x$ in the interval $[0, 1]$. The only way this is possible is if the polynomial is the zero polynomial, meaning all its coefficients must be zero: $a_0 = 0, a_1 = 0, \dots, a_{m-1} = 0$.
	
	Thus, the list $f_0, f_1, \dots, f_{m-1}$ is linearly independent. Because $\V$ contains a linearly independent list of length $m$ for every positive integer $m$, $\V$ cannot be spanned by any finite list of vectors. Therefore, $\V$ is infinite-dimensional.
\end{proof}
\begin{question}
	Suppose $p_0, p_1, \dots, p_m$ are polynomials in $\mathcal{P}_m(\F)$ such that $p_k(2) = 0$ for each $k \in \{0, \dots, m\}$. Prove that $p_0, p_1, \dots, p_m$ is not linearly independent in $\mathcal{P}_m(\F)$.
\end{question}
\begin{proof}
	Because $p_k(2) = 0$ for each $k \in \{0, \dots, m\}$, each polynomial $p_k$ has a root at $2$. Therefore, we can factor out $(z - 2)$ from each polynomial. That is, there exist polynomials $q_0, q_1, \dots, q_m \in \mathcal{P}_{m-1}(\F)$ such that
	\[
	p_k(z) = (z - 2)q_k(z)
	\]
	for each $k \in \{0, \dots, m\}$.
	
	The vector space $\mathcal{P}_{m-1}(\F)$ is spanned by the list of polynomials $1, z, \dots, z^{m-1}$, which has a length of $m$. The list $q_0, q_1, \dots, q_m$ consists of $m+1$ polynomials in $\mathcal{P}_{m-1}(\F)$. 
	
	A fundamental theorem of linear algebra states that in any finite-dimensional vector space, the length of a linearly independent list cannot exceed the length of any spanning list. Because the list $q_0, q_1, \dots, q_m$ has length $m+1$, which is strictly greater than the length of the spanning list $m$, it follows that $q_0, q_1, \dots, q_m$ must be linearly dependent.
	
	Thus, there exist scalars $a_0, a_1, \dots, a_m \in \F$, not all zero, such that
	\[
	a_0 q_0 + a_1 q_1 + \dots + a_m q_m = 0.
	\]
	Multiplying both sides of this equation by $(z - 2)$, we obtain
	\[
	a_0 (z - 2)q_0 + a_1 (z - 2)q_1 + \dots + a_m (z - 2)q_m = 0,
	\]
	which simplifies to
	\[
	a_0 p_0 + a_1 p_1 + \dots + a_m p_m = 0.
	\]
	Because the scalars $a_0, a_1, \dots, a_m$ are not all zero, the list $p_0, p_1, \dots, p_m$ is not linearly independent.
\end{proof}